\documentclass[10pt,a4paper]{article}

\usepackage[top=1.8cm, bottom=1.8cm, left=1.8cm, right=1.8cm, columnsep=0.7cm]{geometry}
\usepackage[spanish]{babel}
\usepackage[utf8]{inputenc}
\usepackage[T1]{fontenc}
\usepackage{multicol}
\usepackage{booktabs}
\usepackage{tabularx}
\usepackage{array}
\usepackage{enumitem}
\usepackage{titlesec}
\usepackage{microtype}
\usepackage{parskip}
\usepackage{xcolor}
\usepackage{fancyhdr}
\usepackage{amsmath}

% ── Colores ────────────────────────────────────────────────────────────────────
\definecolor{uniblue}{RGB}{0,51,102}
\definecolor{lightgray}{RGB}{240,240,240}

% ── Secciones compactas ────────────────────────────────────────────────────────
\titleformat{\section}{\bfseries\small\color{uniblue}}{}{0em}{}[\color{uniblue}\titlerule]
\titleformat{\subsection}{\bfseries\footnotesize}{}{0em}{}
\titlespacing{\section}{0pt}{6pt}{3pt}
\titlespacing{\subsection}{0pt}{4pt}{2pt}

% ── Listas compactas ───────────────────────────────────────────────────────────
\setlist[itemize]{noitemsep, topsep=2pt, leftmargin=1.2em}
\setlist[enumerate]{noitemsep, topsep=2pt, leftmargin=1.4em}

% ── Párrafos ───────────────────────────────────────────────────────────────────
\setlength{\parskip}{3pt}
\setlength{\parindent}{0pt}

% ── Cabecera ───────────────────────────────────────────────────────────────────
\pagestyle{fancy}
\fancyhf{}
\renewcommand{\headrulewidth}{0.4pt}
\fancyhead[L]{\footnotesize\textbf{Manual Técnico — Proyecto Final IELE3338L}}
\fancyhead[R]{\footnotesize Universidad de los Andes · 2026-10}
\fancyfoot[C]{\footnotesize\thepage}

% ── Comandos auxiliares ────────────────────────────────────────────────────────
\newcommand{\topic}[1]{\texttt{\small #1}}
\newcommand{\pkg}[1]{\texttt{\small #1}}

% ══════════════════════════════════════════════════════════════════════════════
\begin{document}

% ── Título ────────────────────────────────────────────────────────────────────
\begin{center}
  {\large\bfseries\color{uniblue}
    Robot Autónomo para Mapeo, Navegación y Manipulación\\[2pt]
    en Pista Logística — Proyecto Final}\\[4pt]
  {\small
    Curso: IELE3338L · Sección 1--2 · Periodo 2026-10\\
    Grupo: \textbf{5} \quad|\quad
    Integrantes: \textbf{Andrés Carlos, Juan Carral}\\
    Profesor: Jorge López Jiménez · AG: Diego Santiago Jiménez Beltrán}
\end{center}
\vspace{-4pt}
\noindent\rule{\linewidth}{0.8pt}
\vspace{2pt}

% ── Inicio de dos columnas ────────────────────────────────────────────────────
\begin{multicols}{2}

% ─────────────────────────────────────────────────────────────────────────────
\section{Descripción del Sistema}

El robot es un móvil diferencial autónomo capaz de construir un mapa de la pista, localizarse, planear rutas, evitar obstáculos y manipular bloques de colores (rojo, azul, verde) para ubicarlos en sus zonas objetivo. El sistema integra cinco subsistemas: percepción, localización, mapeo, navegación y manipulación, coordinados bajo \textbf{ROS2}.

El procesamiento pesado de visión (cámara cenital) se realiza en un PC externo para no saturar la Raspberry Pi, publicando los resultados vía ROS2.

% ─────────────────────────────────────────────────────────────────────────────
\section{Hardware}

\begin{tabular}{@{}lll@{}}
\toprule
\textbf{Componente} & \textbf{Interfaz} & \textbf{Parámetro} \\
\midrule
Raspberry Pi 4B       & —          & Computador principal \\
YDLidar X4 Pro        & Serial     & \texttt{/dev/lidar} \\
IMU BNO055            & I2C bus 1  & Addr \texttt{0x28} \\
Encoders (×2)         & GPIO       & 562 PPR \\
Motores DC (L298N)    & GPIO PWM   & 2 canales \\
Pi Camera 2           & CSI        & 820×640 px \\
ESP32 (brazo)         & Serial USB & \texttt{/dev/ttyUSB0} \\
Arduino (stepper Z)   & Serial USB & \texttt{/dev/arduino} \\
PC externo (visión)   & WebSocket  & Puerto 8765 \\
\bottomrule
\end{tabular}

\vspace{4pt}

\textbf{Parámetros cinemáticos:} rueda $\varnothing = 8\,\text{cm}$, base $L = 20\,\text{cm}$.

GPIO motores: izq.\ 5/6, der.\ 17/27. GPIO encoders: izq.\ 20/21, der.\ 23/24.

% ─────────────────────────────────────────────────────────────────────────────
\section{Arquitectura de Software (ROS2)}

\subsection{Paquetes}

\begin{itemize}
  \item \pkg{sensors} — percepción: IMU, LiDAR, cámara, encoders, brazo, stepper.
  \item \pkg{locomotion} — control de motores (PID) y dashboard de trayectoria.
  \item \pkg{my\_robot\_description} — URDF y configuración EKF.
  \item \pkg{nav2\_config} — SLAM Toolbox + Nav2 + launch completo.
\end{itemize}

\subsection{Nodos principales}

\begin{tabular}{@{}ll@{}}
\toprule
\textbf{Nodo} & \textbf{Publica / Suscribe} \\
\midrule
\texttt{mpu\_node}        & pub \topic{/imu} \\
\texttt{encoders\_node}   & pub \topic{/odom}, \topic{/omega} \\
\texttt{lidar\_node}      & pub \topic{/scan} \\
\texttt{camera\_node}     & pub \topic{/camera/image\_raw} \\
\texttt{vision\_bridge}   & pub \topic{/vision/detections} \\
\texttt{ekf\_node}        & sub \topic{/odom}, \topic{/imu} → pub \topic{/odometry/filtered} \\
\texttt{slam\_toolbox}    & sub \topic{/scan} → pub \topic{/map} \\
\texttt{nav2}             & sub \topic{/map}, \topic{/odometry/filtered} → pub \topic{/cmd\_vel} \\
\texttt{motor\_control}   & sub \topic{/odometry/filtered}, \topic{/ref} → GPIO \\
\texttt{arm\_bridge}      & sub \topic{/arm/mode} → serial ESP32 \\
\texttt{stepper\_node}    & sub \topic{/z\_axis/move} → serial Arduino \\
\bottomrule
\end{tabular}

\subsection{Fusión sensorial}

El filtro EKF (\texttt{robot\_localization}) fusiona la odometría de encoders (\topic{/odom}) con la orientación del BNO055 (\topic{/imu}) a 50\,Hz, publicando \topic{/odometry/filtered} hacia Nav2 y el controlador de motores. La localización cenital (cámara del techo) se integra como fuente adicional de corrección de posición $x, y, \theta$ publicada por el PC externo.

\subsection{Control de movimiento}

El controlador cinemático (\texttt{motor\_control}) implementa un controlador proporcional de posición:
\[
  v = K_s\, e_s, \quad \omega = K_\theta\, e_\theta
\]
con $K_s=0.3$, $K_\theta=2.4$, seguido de un PID de velocidad de rueda ($P=0.775$, $I=28.0$, $D=0.002$) que cierra el lazo sobre los encoders.

% ── Columna derecha ───────────────────────────────────────────────────────────
\columnbreak

% ─────────────────────────────────────────────────────────────────────────────
\section{Secuencia de Arranque}

\begin{enumerate}
  \item Conectar alimentación. Verificar que los dispositivos seriales estén presentes:\\
        \texttt{ls /dev/lidar /dev/ttyUSB0 /dev/arduino}
  \item Sourcear el workspace en la Raspberry Pi:
        \begin{verbatim}
source ~/proyecto_robotica/install/setup.bash
        \end{verbatim}
  \item Lanzar la pila completa (sensores + EKF + SLAM + Nav2 + motores):
        \begin{verbatim}
ros2 launch nav2_config \
  nav2_bringup.launch.py
        \end{verbatim}
  \item En el PC externo (visión cenital), iniciar el cliente WebSocket y esperar conexión con \topic{/vision/detections}.
  \item Lanzar el brazo manipulador:
        \begin{verbatim}
ros2 run sensors arm_bridge_node
ros2 run sensors stepper_node
        \end{verbatim}
  \item Activar el inicio de la misión (pulsador físico o comando de arranque).
  \item Verificar en RViz2 que el mapa, la odometría y el LiDAR estén activos antes de iniciar la ejecución oficial.
\end{enumerate}

% ─────────────────────────────────────────────────────────────────────────────
\section{Herramientas de Monitoreo}

\begin{itemize}
  \item \textbf{Dashboard de trayectoria:} \texttt{ros2 run locomotion graphical\_odom} → \texttt{http://<ip>:5000} (Plotly, UDP).
  \item \textbf{RViz2:} visualización de \topic{/map}, \topic{/scan}, pose del robot y ruta planeada.
  \item \textbf{Estado del brazo:} \texttt{ros2 topic echo /arm/status}.
  \item \textbf{Posición Z:} \texttt{ros2 topic echo /z\_axis/position}.
\end{itemize}

% ─────────────────────────────────────────────────────────────────────────────
\section{Manipulación de Bloques}

El brazo manipulador (4 servos, control mediante ESP32) opera en cuatro modos publicados en \topic{/arm/mode}: \texttt{live} (teleoperación en tiempo real), \texttt{record} (grabación a CSV), \texttt{replay} (reproducción de secuencia grabada) y \texttt{stop} (posición home $= [90°, 90°, 90°, 90°]$).

El eje Z (motor paso a paso controlado por Arduino vía AccelStepper) permite movimiento relativo en mm publicando en \topic{/z\_axis/move}. Posición actual disponible en \topic{/z\_axis/position}.

La detección de bloques por color (rojo/azul/verde, cubos de 15\,cm) se realiza en el PC externo usando segmentación HSV sobre la imagen cenital; los resultados (posición, color, estado) se publican en \topic{/vision/detections} (JSON).

% ─────────────────────────────────────────────────────────────────────────────
\section{Procedimientos de Seguridad}

\begin{itemize}
  \item \textbf{Parada de emergencia:} \texttt{Ctrl+C} en el launch detiene todos los nodos y envía la señal \texttt{stop} al brazo antes de cerrar el serial.
  \item \textbf{Límite de velocidad:} la salida PWM de los motores se limita a $[-1.0, 1.0]$ en el PID para evitar sobrecorriente.
  \item \textbf{Detención por llegada:} el controlador se detiene cuando $e_s < 5\,\text{cm}$.
  \item \textbf{Gálibo y puente:} el sistema de visión debe marcar como celdas no transitables todas las regiones del mapa que excedan la altura del gálibo (30\,cm) cuando no corresponda atravesarlo.
  \item \textbf{Condición insegura:} si se detecta pérdida de comunicación serial o falla del BNO055, el nodo correspondiente emite \texttt{WARN} pero no detiene el sistema; el EKF puede continuar solo con odometría.
\end{itemize}

% ─────────────────────────────────────────────────────────────────────────────
\section{Limitaciones Conocidas}

\begin{itemize}
  \item La deriva de odometría se incrementa en superficies irregulares; la corrección cenital es esencial para misiones largas.
  \item La cámara del robot (Pi Camera) no está utilizada directamente para navegación, solo para el streaming al PC externo.
  \item El nodo \texttt{motor\_command} (\topic{/motor\_cmd}) y \texttt{motor\_control} (GPIO directo) no deben ejecutarse simultáneamente (conflicto de pines GPIO).
\end{itemize}

\end{multicols}

\end{document}
